\section{Method}
\label{sec:method}

\begin{figure}[t]
  \centering
  \includegraphics[width=\columnwidth]{figures/figure_1_pipeline.png}
  \caption{
    Pipeline overview. Phase~1 constructs scene-specific lever candidates
    from an editor-feasible intervention vocabulary. Phase~2 realises them via bounded stochastic
    generation with validity auditing. Phase~3 scores accepted edits with
    an auxiliary perception model.
  }
  \label{fig:pipeline}
\end{figure}

\subsection{Problem Setup}

Let $x \in \mathcal{X}$ denote a street-view image and
$a \in \mathcal{A}$ a perceptual attribute (e.g.\ safety). We take the
explanatory atom to be a \emph{lever intervention}
\begin{equation}
  l = (c, s, d, \tau),
\end{equation}
where $c$ is a lever concept (e.g.\ greenery, graffiti), $s$ is the
grounded scene support, $d$ the intervention direction (add, remove,
repair), and $\tau$ a constrained edit template.

For image $x$, let $L(x)$ denote the scene-specific candidate set
instantiated from an editor-feasible intervention vocabulary. Our method identifies the
subset that can be realised as valid edits:
\begin{equation}
  V(x) = \{ l_i \in L(x) \mid \exists\, \tilde{x}_i \text{ passing the validity audit} \}.
\end{equation}
Each lever is tested independently against the unedited original; edits
are never composed sequentially.

\subsection{Phase 1: Lever Candidate Construction}

We define an \emph{editor-feasible intervention vocabulary}: a bounded
set of lever concepts that each satisfy three inclusion criteria.
\textbf{(i)}~\emph{Empirical grounding}: the cue has prior evidence
linking it to perceived safety in the urban perception
literature~\cite{jacobs1961,newman1972,li2015greenery,jiang2018brokenwindows,portnov2020lighting}.
\textbf{(ii)}~\emph{Semantic distinctness}: each item targets a
non-overlapping scene element so that lever attributions remain
interpretable (e.g.\ \emph{litter removal} and \emph{surface cleaning}
act on different substrates).
\textbf{(iii)}~\emph{Prompt-only editability}: the intervention must be
realisable by a current diffusion editor without segmentation masks,
which excludes structurally complex changes such as road-geometry
redesign or building-mass alteration.

The twelve resulting concepts are organised into four families
reflecting distinct urban-design mechanisms
(Table~\ref{tab:vocabulary}; full list in
Appendix~\ref{sec:intervention_vocabulary}):
\textbf{Physical Maintenance} covers disorder removal and surface
repair, grounded in broken-windows and perceived-upkeep
theory~\cite{newman1972,jiang2018brokenwindows};
\textbf{Environmental Amenity} covers greenery and lighting additions,
grounded in biophilic safety and nighttime-visibility
research~\cite{li2015greenery,portnov2020lighting};
\textbf{Visual Legibility} covers signage and storefront transparency,
grounded in Jacobs' active-edges principle~\cite{jacobs1961};
and \textbf{Mobility Infrastructure} covers road-surface markings,
grounded in pedestrian-safety and walkability
cues~\cite{quintana2025}.
For each image, a VLM planner grounds vocabulary concepts to visible
scene locations, yielding candidates such as \emph{remove graffiti from
this wall} or \emph{add modest greenery near this entrance}.

\begin{table}[t]
  \centering
  \caption{Editor-feasible intervention vocabulary. Each concept
  satisfies the empirical-grounding, semantic-distinctness, and
  prompt-only-editability criteria above. Theory anchors indicate
  the primary literature motivation per family.}
  \label{tab:vocabulary}
  \scriptsize
  \begin{tabular}{p{0.22\linewidth}p{0.50\linewidth}l}
    \toprule
    Family & Lever concepts & Anchor \\
    \midrule
    Physical Maint.\ &
      graffiti removal, litter removal,
      facade repair, surface cleaning,
      shutter repair &
      \cite{newman1972,jiang2018brokenwindows} \\
    Environ.\ Amenity &
      localised greenery addition,
      lighting repair,
      tree canopy management &
      \cite{li2015greenery,portnov2020lighting} \\
    Visual Legibility &
      signage decluttering,
      storefront transparency increase &
      \cite{jacobs1961} \\
    Mobility Infra.\ &
      crosswalk repainting,
      lane marking repainting &
      \cite{quintana2025} \\
    \bottomrule
  \end{tabular}
\end{table}

\subsection{Phase 2: Bounded Stochastic Intervention}

For each candidate $l_i$, a prompt-only image generator $G$ produces
\begin{equation}
  x'_{i,j} = G(x;\, c_i, s_i, d_i, \tau_i, \varphi_j),
\end{equation}
where $\varphi_j$ indexes stochastic draws under a bounded retry budget
$T$. Currently $G$ is \texttt{flux-kontext-pro}; the framework is
generator-agnostic.

Each candidate is subjected to an automated validity audit performed by
a vision--language model (GPT-5.2 in the current implementation),
evaluating: \textbf{(1)}~same-place preservation,
\textbf{(2)}~locality to the intended support, \textbf{(3)}~realism,
and \textbf{(4)}~plausibility of the intended lever. If at least one
draw passes within $T$ attempts, the first accepted edit $\tilde{x}_i$
is retained and $l_i$ enters $V(x)$. We report valid-edit coverage as
$\mathrm{Coverage}(x) = |V(x)| / |L(x)|$.

\subsection{Phase 3: Auxiliary Scoring}

For each accepted edit $\tilde{x}_i$, we compute
\begin{equation}
  \Delta_i = f_a(\tilde{x}_i) - f_a(x),
\end{equation}
where $f_a$ is a ViT-B/16 perception model pretrained on the MIT Place
Pulse~2.0 pairwise safety dataset~\cite{hou2024globalstreetscapes},
which outputs a continuous score on a 0--10 scale (higher = safer).
This is strictly an auxiliary proxy, not the main estimand; $f_a$ has
not been validated on edited images, so $\Delta_i$ should be treated as
a ranking signal rather than a calibrated effect size. The auxiliary-promising set is
\begin{equation}
  E_{\mathrm{aux}}(x,a) = \{ l_i \in V(x) \mid \Delta_i >
  \theta_{\mathrm{aux}} \},
\end{equation}
where $\theta_{\mathrm{aux}}$ is an exploratory threshold.
Human-grounded effectiveness via pairwise evaluation remains future
work.%
\footnote{Our method is constrained interventional search, not
pixel-level decomposition or causal identification. ``Multiple levers''
means distinct interventions that each pass validity checks
independently; it does not imply additive effects. Locality is audited,
not pixel-guaranteed.}
